% Part: lambda-calculus
% Chapter: church-rosser
% Section: beta-reduction

\documentclass[../../../include/open-logic-section]{subfiles}

\begin{document}

\olfileid{lam}{cr}{b}

\olsection{$\beta$-হ্রাস}

\begin{lem}\ollabel{lem:one-par}
  $M \bredone M'$ হলে $M \bredpar M'$।
\end{lem}
\begin{proof} $M \bredone M'$-এর ব্যুৎপত্তির উপর আরোহ প্রয়োগ
  করি। শেষ ধাপটি বিমূর্তন বা প্রয়োগের সামঞ্জস্য-বিধিগুলির একটি হলে,
  আরোহ-অনুমান এবং সমান্তরাল হ্রাসের সংশ্লিষ্ট বিধি থেকে সিদ্ধান্তটি আসে।
% BN-SRC-295 TARGET-CORRECTION-BEGIN
\par\smallskip\noindent\textnormal{\small\textbf{উৎস-সংশোধন (BN-SRC-295):} হিমায়িত প্রমাণটি $M\bredone M'$ হলেই $M$-কে মূলস্থ রিডেক্স $(\lambd[x][N])Q$ বলে ধরে নেয়। কিন্তু $\bredone$ ক্ষুদ্রতম সামঞ্জস্যপূর্ণ সম্বন্ধ হওয়ায় সংকোচন বিমূর্তনের ভিতরে বা প্রয়োগের যেকোনো উপপদেও ঘটতে পারে। বাংলা প্রমাণে তাই !!{derivation}-এর উপর আরোহ যোগ করা হয়েছে: তিনটি সামঞ্জস্য-ক্ষেত্রে আরোহ-অনুমান ও সমান্তরাল বিধি ব্যবহৃত হয়, আর নিচের মূলস্থ ক্ষেত্রটি হিমায়িত যুক্তিটিই।} % BN-SRC-295 TARGET-CORRECTION-END
  শেষ ধাপটি মূল সংকোচন হলে $M$ হলো
  $(\lambd[x][N])Q$, $M'$ হলো $\Subst{N}{Q}{x}$, কোনো
  $x$, $N$ ও~$Q$-এর জন্য। \olref[pb]{thm:refl} অনুসারে
  $N \bredpar N$ এবং $Q \bredpar Q$; অতএব
  \olref[pb]{defn:bredpar}\olref[pb]{defn:bredpar4} অনুসারে অবিলম্বে
  $(\lambd[x][N])Q \bredpar \Subst{N}{Q}{x}$।
\end{proof}

\begin{lem}\ollabel{lem:par-red}
  $M \bredpar M'$ হলে $M \bred M'$।
\end{lem}

\begin{proof} $M \bredpar M'$-এর !!{derivation}-এর উপর আরোহ প্রয়োগ
  করি।
  \begin{enumerate}
  \item শেষ বিধিটি \olref[pb]{defn:bredpar1}: তখন $M$ ও~$M'$ উভয়েই
    কেবল~$x$, এবং $x \bred x$।
  \item শেষ বিধিটি \olref[pb]{defn:bredpar2}: কোনো $x$, $N$, $N'$-এর
    জন্য $M$ হলো $\lambd[x][N]$ এবং $M'$ হলো $\lambd[x][N']$, যেখানে
    $N \bredpar N'$। আরোহ-অনুমান থেকে $N \bred N'$। তখন $N \bred N'$
    যে $\bredone$-সংকোচনধারা দিয়ে হয়, সেই একই ধারা থেকে
    $\lambd[x][N] \bred \lambd[x][N']$।
  \item শেষ বিধিটি \olref[pb]{defn:bredpar3}: কোনো $P$, $Q$, $P'$,
    $Q'$-এর জন্য $M$ হলো $PQ$ এবং $M'$ হলো $P'Q'$, যেখানে
    $P \bredpar P'$ ও~$Q \bredpar Q'$। আরোহ-অনুমান থেকে
    $P \bred P'$ ও~$Q \bred Q'$। অতএব প্রথমে $P \bred P'$ এবং তারপর
    $Q \bred Q'$ প্রয়োগের হ্রাসধারা দিয়ে $PQ \bred P'Q'$।
  \item শেষ বিধিটি \olref[pb]{defn:bredpar4}: কোনো $x$, $N$, $N'$,
    $Q$, $Q'$-এর জন্য $M$ হলো $(\lambd[x][N])Q$ এবং $M'$ হলো
    $\Subst{N'}{Q'}{x}$, যেখানে $N \bredpar N'$ ও~$Q \bredpar Q'$।
% BN-SRC-296 TARGET-CORRECTION-BEGIN
\par\smallskip\noindent\textnormal{\small\textbf{উৎস-সংশোধন (BN-SRC-296):} হিমায়িত চলকতালিকায় $N$, $M'$, $Q$, $Q'$ লেখা, যদিও পূর্বধারণা ও পরের প্রতিটি লাইনে দরকার $N$, $N'$, $Q$, $Q'$। বাংলা তালিকায় অনুপস্থিত $N'$ পুনরুদ্ধার করা হয়েছে।} % BN-SRC-296 TARGET-CORRECTION-END
    আরোহ-অনুমান থেকে $Q \bred Q'$ ও~$N \bred N'$। তাই প্রথমে
    $N \bred N'$, তারপর $Q \bred Q'$, এবং শেষে
    $(\lambd[x][N'])Q'$-কে $\Subst{N'}{Q'}{x}$-এ সংকুচিত করে
    $(\lambd[x][N])Q \bred \Subst{N'}{Q'}{x}$ পাই।
  \end{enumerate}
\end{proof}

\begin{lem}\ollabel{lem:str}
  $\bredpar$-কে ধারণকারী ক্ষুদ্রতম পরিযায়ী সম্বন্ধ হলো~$\bred$।
\end{lem}

\begin{proof}
  $\xred$-কে $\bredpar$-কে ধারণকারী ক্ষুদ্রতম পরিযায়ী সম্বন্ধ ধরি।

  $\bred \subseteq \xred$: ধরি $M \bred M'$; অর্থাৎ $M \ident M_1
  \bredone \dots \bredone M_k \ident M'$। \olref{lem:one-par} অনুসারে
  $M \ident M_1 \bredpar \dots \bredpar M_k \ident M'$। যেহেতু
  $\xred$ সম্বন্ধটি $\bredpar$-কে ধারণ করে এবং পরিযায়ী, তাই
  $M \xred M'$।

  $\xred \subseteq \bred$: ধরি $M \xred M'$; অর্থাৎ $M
  \ident M_1 \bredpar \dots \bredpar M_k \ident M'$।
  \olref{lem:par-red} অনুসারে $M \ident M_1
  \bred \dots \bred M_k \ident M'$। যেহেতু $\bred$ পরিযায়ী, তাই
  $M \bred M'$।
\end{proof}

\begin{thm}\ollabel{thm:cr}
  $\bred$ চার্চ--রসার ধর্ম পূরণ করে।
\end{thm}

\begin{proof}
  \olref[dap]{thm:str}, \olref[pb]{thm:cr} ও~\olref{lem:str} থেকে
  অবিলম্বে।
\end{proof}

\end{document}
